\documentclass[12pt,a4paper]{article}
\usepackage{amsmath}
\usepackage{amssymb}
\usepackage{amsthm}
\usepackage{geometry}
\usepackage{xcolor}
\usepackage{hyperref}

\geometry{margin=1in}

% Define theorem environments
\theoremstyle{definition}
\newtheorem{definition}{Definition}[section]
\newtheorem{theorem}{Theorem}[section]
\newtheorem{proposition}{Proposition}[section]
\newtheorem{lemma}{Lemma}[section]

\title{COMP0083: Kernel Methods}
\author{Lecture Notes}
\date{}

\begin{document}

\maketitle
\tableofcontents
\newpage

\section{Introduction to Kernels}

\begin{enumerate}
    \item Kernels let us compare objects on the basis of features
    \item Kernel methods can control smoothness and avoid overfitting/underfitting
\end{enumerate}

\section{Inner Product}

Let $H$ be a vector space over $\mathbb{R}$. A function $\langle \cdot, \cdot \rangle_H : H \times H \to \mathbb{R}$ is an inner product if:

\begin{enumerate}
    \item[(i)] \textbf{Linear:} $\langle \alpha f, \beta g \rangle_H = \alpha \langle f, g \rangle_H + \alpha_2 \langle f_2, g \rangle_H$
    \item[(ii)] \textbf{Symmetric:} $\langle f, g \rangle_H = \langle g, f \rangle_H$
    \item[(iii)] $\langle f, f \rangle_H \geq 0$ and $\langle f, f \rangle_H = 0 \Leftrightarrow f = 0$ (positive definiteness)
\end{enumerate}

\[
\|f\|_H := \sqrt{\langle f, f \rangle_H}
\]

\section{Inner Product on $\mathbb{R}^n$}

A function $\langle \cdot, \cdot \rangle : \mathbb{R}^n \times \mathbb{R}^n \to \mathbb{R}$ is an inner product on $\mathbb{R}^n$ if:

\[
\exists \text{ symmetric positive definite matrix } M \text{ s.t. } \langle x, y \rangle = x^T M y \text{ for } \forall x, y \in \mathbb{R}^n
\]

For our product with $M = I$:

\[
\langle x, y \rangle = x^T y
\]

\section{Hilbert Space}

A Hilbert space is an inner product space containing all Cauchy sequence limits.

\section{Kernel Definition}

Let $X$ be a non-empty set. A function $k : X \times X \to \mathbb{R}$ is a kernel if there exists a Hilbert space $H$ and a feature map $\phi : X \to H$ such that:

\[
k(x, x') = \langle \phi(x), \phi(x') \rangle_H
\]

\section{Key Properties of Kernels}

\begin{enumerate}
    \item A single kernel can correspond to several possible features.
    
    \item \textbf{Sums of kernels are kernels:} Given $\alpha > 0$, $k, k_2, k_3$ are kernels on $X$, then $\alpha k$ and $k_1 + k_2$ are kernels on $X$.
    
    \item A difference of kernels may not be a kernel.
    
    \item \textbf{Mappings between spaces:} Let $X$ and $\tilde{X}$ be sets and a map $A : X \to \tilde{X}$. Define kernel on $\tilde{X}$. Then $k(A(x), A(x'))$ is a kernel on $X$.
    
    \item \textbf{Products of kernels are kernels:} Given $k_1$ on $X_1$ and $k_2$ on $X_2$, then $k_1 \times k_2$ is a kernel on $X_1 \times X_2$. If $X_1 = k_1 = X$, then $k := k_1 \times k_2$ is a kernel on $X$.
\end{enumerate}

\section{Frobenius Inner Product}

\[
\langle A, B \rangle_F = \sum_{ij} A_{ij} B_{ij} = \text{tr}(A^T B), \quad A^* \text{ is conjugate-transpose of } A
\]

\section{Polynomial Kernel}

\[
k(x, x') := (\langle x, x' \rangle + c)^m
\]

where $m$ must be an integer and $m \geq 1$, $c \geq 0$.

\section{Convergence Radius}

For $r \in (0, \infty]$, with $a_n \geq 0$ for all $n \geq 0$:

\[
f(z) = \sum_{n=0}^{\infty} a_n z^n, \quad |z| < r, \quad z \in \mathbb{R}
\]

Define $X_0$ to be the $\sqrt{r}$-ball in $\mathbb{R}^d$, so $\|x\| < \sqrt{r}$:

\[
k(x, x') = f(\langle x, x' \rangle) = \sum_{n=0}^{\infty} a_n \langle x, x' \rangle^n
\]

\section{Sequence Spaces}

The space $\ell_2$ comprises all sequences $a := \{a_i\}_{i>1}$ for which:

\[
\|a\|_{\ell_2}^2 = \sum_{i=1}^{\infty} a_i^2 < \infty
\]

\newpage

\section{Properties of Kernels and Positive Definiteness}

\begin{enumerate}
    \item Kernels are positive definite. Reverse also holds.
\end{enumerate}

\subsection{Norm}

Let $F$ be a vector space over the field $\mathbb{R}$ (or $\mathbb{C}$). A function $\|\cdot\|_F : F \to [0, \infty)$ is said to be a norm if:

\begin{enumerate}
    \item[(i)] $\|f\|_F = 0 \Leftrightarrow f = 0$ (norm separates points)
    \item[(ii)] $\|\lambda f\|_F = |\lambda| \|f\|_F$, $\forall \lambda \in \mathbb{R}$, $\forall f \in F$ (positive homogeneity)
    \item[(iii)] $\|f + g\|_F \leq \|f\|_F + \|g\|_F$, $\forall f, g \in F$ (triangle inequality)
\end{enumerate}

In every normed vector space, one can define a metric induced by the norm: $d(f, g) = \|f - g\|_F$.

In every inner product vector space, one can define a norm induced by the inner product:

\[
\|f\|_F = \langle f, f \rangle_F^{1/2}
\]

\subsection{Orthogonal Complement}

The orthogonal complement of $M$: $M^{\perp} = \{g \in F : \langle f, g \rangle = 0, \forall f \in M\}$

$M^{\perp}$ is a linear subspace of $F$; $M \cap M^{\perp} = \{0\}$.

\subsection{Key Relations in Inner Product Space}

\begin{enumerate}
    \item \textbf{Cauchy-Schwarz inequality:} $|\langle f, g \rangle| \leq \|f\|_F \|g\|_F$
    
    \item \textbf{Parallelogram law:} $2\|f\|_F^2 + 2\|g\|_F^2 = \|f + g\|_F^2 + \|f - g\|_F^2$
    
    \item \textbf{Polarization identity:} $4\langle f, g \rangle = \|f + g\|_F^2 - \|f - g\|_F^2$
    
    \item \textbf{Pythagorean theorem:} $f \perp g \Rightarrow \|f\|_F^2 + \|g\|_F^2 = \|f + g\|_F^2$
\end{enumerate}

\subsection{Convergence and Cauchy Sequences}

From $\|f_n - f_m\|_F \leq \|f_n - f\|_F + \|f - f_m\|_F$, we get: Convergent $\Rightarrow$ Cauchy, Cauchy $\Rightarrow$ Convergent?

Example: $1, 1.4, 1.41, 1.414, 1.4142, \ldots$ is Cauchy in $\mathbb{Q}$ which does not converge because $\sqrt{2} \notin \mathbb{Q}$.

\subsection{Complete Metric Spaces}

A metric space $F$ is said to be complete if every Cauchy sequence $\{f_n\}_{n=1}^{\infty}$ in $F$ converges: it has limit and this limit is in $F$, i.e., one can find $f \in F$ s.t. $\lim \|f_n - f\|_F = 0$.

\begin{center}
\textbf{Complete + Norm = Banach} \quad \textbf{Complete + Inner product = Hilbert}
\end{center}

\subsection{Closed vs. Complete}

\textbf{Closed v.s. Complete:} Closed: $M \subseteq F$ is closed in $F$ $\Leftrightarrow$ it contains limits of all sequences in $M$ that converge in $F$.

Complete: $M$ is complete if all Cauchy sequence in $M$ converge in $M$.

\subsection{Linear Operators}

\textbf{Linear operator:} $A: F \to G$, where $F$ and $G$ are both vector spaces over $\mathbb{R}$. $A$ is said to be a linear operator if:

\[
A(\alpha_1 f_1 + \alpha_2 f_2) = \alpha_1 A(f_1) + \alpha_2 A(f_2), \quad \forall \alpha_1, \alpha_2 \in \mathbb{R}, \quad f_1, f_2 \in F
\]

Operators with $G = \mathbb{R}$ are called \textbf{functionals}.

\subsection{Operator Norm}

\textbf{Operator norm:} The operator norm of a linear operator $A : F \to G$ is defined as:

\[
\|A\| = \sup_{f \in F} \frac{\|A f\|_G}{\|f\|_F}
\]

If $\|A\| < \infty$, $A$ is called a \textbf{bounded linear operator}. $\|A\|$ is the smallest number $\lambda$ s.t. the inequality $\|A f\|_G \leq \lambda \|f\|_F$ holds for all $f \in F$.

\subsection{Continuity}

\textbf{Continuity:} $A: F \to G$, where $F$ and $G$ are both normed vector spaces over $\mathbb{R}$. $A$ is said to be continuous at $f_0 \in F$, if $\forall \varepsilon \geq 0$, $\exists \delta = S(\varepsilon, f_0) > 0$, s.t.

\[
\|f - f_0\|_F < \delta \Rightarrow \|A f - A f_0\|_G < \varepsilon
\]

$A$ is said to be continuous on $F$ if it is continuous at every point of $F$.

Example: For $g \in F$, $A_g : F \to \mathbb{R}$ defined with $A_g(f) = \langle f, g \rangle_F$ is continuous on $F$.

$|A_g f_1 - A_g f_2| = |\langle f_1 - f_2, g \rangle| \leq \|g\|_F \|f_1 - f_2\|_F$, so take $\delta = \frac{\varepsilon}{\|g\|_F}$ (also Lipschitz?)

\newpage

\section{Continuous and Bounded Operators}

\begin{enumerate}
    \item \textbf{Continuous function:} $A: F \to G$ maps to open (closed) sets from open (closed) sets.
    
    \item \textbf{Bounded linear operator:} $A : F \to G$ maps bounded sets to bounded sets.
    
    \item \textbf{Continuous operator = Bounded operator:} Let $(F, \|\cdot\|_F)$ and $(G, \|\cdot\|_G)$ be normed linear spaces. If $L$ is a linear operator, then the following three conditions are equivalent:
    \begin{enumerate}
        \item[(i)] $L$ is bounded operator
        \item[(ii)] $L$ is continuous on $F$
        \item[(iii)] $L$ is continuous at one point of $F$
    \end{enumerate}
\end{enumerate}

\section{Riesz Representation Theorem}

Let $H$ be a Hilbert space and let $f : H \to \mathbb{F}$ be a bounded linear functional. Then $\exists! y \in H$ s.t.

\[
f(x) = \langle x, y \rangle_H \quad \forall x \in H \quad \text{and} \quad \|f\| = \|y\|_H
\]

Every continuous linear score is just a dot-product with some hidden vector.

$f \in H'$ (continuous dual of $H$), $\tilde{F} := \{f : H \to \mathbb{F} \mid f \text{ is linear and } \|f\|_{H'} = \sup\frac{\|f\|}{\|x\|_H}\}$.

The Riesz isomorphism is the mapping $J: H \to H'$, $J(y)(x) = \langle x, y \rangle_H$, and it is bijective and norm-preserving ($\|f\|_H = \|y\|_H$), so it is isometric isomorphism.

\section{Evaluation Functional}

Given a space of functions $X$ and a fixed point $t_0$ in the domain, the evaluation functional at $t_0$ : $E_{t_0} : X \to \mathbb{F}$, $E_{t_0}(f) = f(t_0)$.

It is linear but not always bounded.

\section{Reproducing Kernel Hilbert Space}

Reproducing kernel Hilbert Space $\Leftrightarrow$ every evaluation functional $E_x: H \to \mathbb{F}$, $E_x(f) = f(x)$, is bounded with respect to norm of $H$.

Bounded evaluations $\Rightarrow$ Reproducing kernel: Because $E_x \in H'$ is bounded, from Riesz, $\exists k(\cdot, x) = k_x \in H$ s.t. $f(x) = E_x(f) = \langle f, k(\cdot, x) \rangle_H$ $\forall f \in H$.

The function $k : X \times X \to \mathbb{F}$, $k(x, y) := k_y(x) = \langle k_y, k_x \rangle_H$, is called \textbf{reproducing kernel}.

\textbf{Property:} $f(x) = \langle f, k(\cdot, x) \rangle_H$ is the \textbf{reproducing property}.

Reproducing kernel $\Rightarrow$ bounded evaluation.

\subsection{Moore-Aronszajn Theorem}

If $k$ is positive-definite on $X \times X$, $\exists!$ RKHS $H_k$ where kernel is $k$ s.t.

\[
f(x) = \langle f, k(\cdot, x) \rangle_{H_k}, \quad \forall f \in H_k, \quad x \in X
\]

$|f(x)| \leq \|f\|_{H_k} \|k(\cdot, x)\|_{H_k}$ and hence $E_x$ is bounded.

\section{Properties of Reproducing Kernels}

\begin{enumerate}
    \item \textbf{Reproducing kernel is unique and RKHS:} It is a reproducing kernel Hilbert space $\Leftrightarrow$ it has reproducing kernel.
    
    \item \textbf{Every reproducing kernel is a kernel:} We can take $\phi : x \mapsto k(\cdot, x)$, $k(x, y) = \langle k(\cdot, x), k(\cdot, y) \rangle_{H_k}$, i.e., $\mathbb{RHS} H$ is a feature space.
\end{enumerate}

\newpage

\section{Orthonormal Basis}

In Hilbert space $H$, an orthonormal basis is a family $\{e_k\}_{k \in B}$ of elements of $H$ satisfying:

\begin{enumerate}
    \item[(i)] \textbf{Orthogonality:} $\langle e_i, e_j \rangle = 0$ $\forall i, j \in B$ with $i \neq j$
    \item[(ii)] \textbf{Normalization:} $\|e_k\|_H = 1$ for $b \in B$
    \item[(iii)] \textbf{Completeness:} linear span of family $\{e_k, k \in B\}$ is dense in $H$
\end{enumerate}

\section{Spectral Theorem}

Let $X$ be a compact metric space and $k : X \times X \to \mathbb{R}$ a continuous kernel. Define

\[
H = \left\{ f = \sum_{j \in J} a_j e_j : \left\{ \frac{a_j}{\sqrt{\lambda_j}} \in \ell^2(J) \right\} \right\} \quad \text{with inner product:}
\]

\[
\left\langle \sum_{j \in J} a_j e_j, \sum_{b \in J} b_j e_j \right\rangle_H = \sum_{j \in J} \frac{a_j b_j}{\lambda_j}
\]

Then $H = H_k$.

\section{Spectral Decomposition}

Let $F$ be a Hilbert space, and $T : F \to F$ a compact, self-adjoint operator. There is an at most countable orthonormal set (CONS) $\{e_j\}_{j \in J}$ of $F$ and $\{\lambda_j\}_{j \in J}$ with $|\lambda_1| > |\lambda_2| > \ldots > 0$ converging to $0$, such that

\[
f = \sum_{j \in J} \langle f, e_j \rangle e_j \quad \text{and} \quad T f = \sum_{j \in J} \lambda_j \langle f, e_j \rangle e_j, \quad f \in F
\]

\section{Integral Linear Operator}

\[
(K f)(x) = \int_X f(y) K(x, y) \, dy
\]

It is linear because $[K(\alpha f_1 + \beta f_2)](x) = \alpha [Kf_1](x) + \beta [Kf_2](x)$.

\section{Mercer's Theorem}

Let $k$ be a continuous kernel on compact metric space $X$, and let $\nu$ be a finite Borel measure on $X$ with $\text{supp}(\nu) = X$. Then $\forall x, y \in X$:

\[
k(x, y) = \sum_{j \in J} \lambda_j e_j(x) e_j(y)
\]

\{$\lambda_j e_j$\} and $\{e_j\}_{j \in J}$ are eigenvalues and eigenfunctions of $T_{k,\nu}$ on the RKHS, and the convergence of sum is uniform on $X \times X$ and absolute for each pair $(x, y) \in X \times X$.

\[
k(x, y) = \sum_{j \in J} \lambda_j e_j(x) e_j(y) = \langle \sqrt{\lambda_j} e_j(x), \sqrt{\lambda_j} e_j(y) \rangle_u
\]

\section{Alternative Formulation}

Item 43 can be viewed in a different way:

Rescale the eigenfunctions: $\varphi_j := \sqrt{\lambda_j} e_j$ and $\langle \varphi_i, \varphi_j \rangle_H = \delta_{ij}$, so $\langle e_i, e_j \rangle_H = \frac{\delta_{ij}}{\lambda_j}$.

$M$ is Gram matrix of the original eigenfunctions $\{e_j\}$ in the RKHS.

So any function $f(x) = \sum f_n e_n(x) = \sum \frac{\hat{f}_n}{\sqrt{\lambda_n}} \varphi_n(x)$

\[
\|f\|_H^2 = \sum \frac{|\hat{f}_n|^2}{\lambda_n}
\]

For the evaluation functional $E_x(f) = f(x)$, Cauchy-Schwarz gives:

\[
|f(x)|^2 = \left| E_x(f) \right|^2 \leq \left( \sum \frac{\hat{f}_n^2}{\lambda_n} \right) \left( \sum \|e_n\|_H^2 k(x, x) \right) = \|f\|_H^2 k(x, x)
\]

For $E_x$ bounded, we need $\left\{ \frac{|f_n|}{\sqrt{\lambda_n}} \right\} \in \ell^2(J)$ and $k(x, x) < \infty$.

For other ordinary dot product in $L_2$, $\lambda_n = 1 \Rightarrow$ the series diverges $\Rightarrow$ unbounded $\Rightarrow$ not RKHS.

\section{Infinite Feature Space via Fourier Series}

\[
k(x, y) = \sum_{j \in J} \lambda_j e_j(x) e_j(y) = \sum_{j \in J} \lambda_j \exp(i \omega_j (x - y))
\]

which is translation invariant.

\section{Eigenvalues and Eigenfunctions}

Since $\{\lambda_j\}_{j \in J}$ and $\{e_j\}_{j \in J}$ are eigenvalues and eigenfunctions of $T_k$, then we have:

\[
T_k e_j(x) = \int e_j(y) k(x, y) \, \nu(y) dy = \int e_j(y) \sum_{j'} \lambda_{j'} e_{j'}(x) e_{j'}(y) \, \nu(y) dy = \lambda_j e_j(x)
\]

\newpage

\section{Small RKHS Norm and Smooth Functions}

Small RKHS norm results in smooth functions.

\subsection{Two Functions Identical in RKHS Norm}

Two functions identical in RKHS norm agree at every point:

\[
|f(x) - g(x)| = |E_x(f - g)| = |\langle f - g, k_x \rangle_H| \leq \|f - g\|_H \|k_x\|_H = 0
\]

\[
\|k_x\|_H = \langle k(\cdot, x), k(\cdot, x) \rangle_H^{\frac{1}{2}} = k(x, x)^{\frac{1}{2}}
\]

\subsection{Distance Between Means in Feature Space}

\[
\left\| \frac{1}{m} \sum_{i=1}^m \phi(x_i) - \frac{1}{n} \sum_{j=1}^n \phi(y_j) \right\|_H^2 \quad \to \, \text{distance between their means in feature space}
\]

When $\phi(x) = x$, distinguish means. When $\phi(x) = Lx x^2$, distinguish means and variances.

\section{Principal Component Analysis (PCA)}

PCA: First principal component (max. variance)

\[
u_1 = \arg\max_{\|u\|=1} \frac{1}{n} \sum_{i=1}^n \left( u^T \left( x_i - \frac{1}{n} \sum_{j=1}^n x_j \right) \right)^2 = \arg\max_{\|u\|=1} u^T C u
\]

where

\[
C = \frac{1}{n} \sum_{i=1}^n \left( x_i - \frac{1}{n} \sum_{j=1}^n x_j \right) \left( x_i - \frac{1}{n} \sum_{j=1}^n x_j \right)^T = \frac{1}{n} X H X^T, \quad X = [x_1 \ldots x_n]^T
\]

\[
H = I - \frac{1}{n} \mathbf{1}_{n \times n} \quad \text{where } \mathbf{1}_{n \times n} \text{ is a matrix of ones}
\]

If $C = Q \Lambda Q^T$, then $u^T C u = (Q^T u)^T \Lambda (Q^T u) = \sum_i \lambda_i \tilde{u}_i^2 \leq \lambda_{\max} \sum_i \tilde{u}_i^2$

where $\tilde{u} = Q^T u$ are the coordinates of $u$ in the eigenbasis. If $\|u\|^2 = 1$, then $\|\tilde{u}\|^2 = u^T Q Q^T u = u^T u = 1$.

\section{Kernel PCA}

\[
f_1 = \arg\max_{\|f\|_H \leq 1} \frac{1}{n} \sum_{i=1}^n \left( \langle f, \phi(x_i) \rangle_H - \frac{1}{n} \sum_{j=1}^n \langle f, \phi(x_j) \rangle_H \right)^2
\]

\[
= \arg\max_{\|f\|_H \leq 1} \frac{1}{n} \sum_{i=1}^n \left( f(x_i) - \frac{1}{n} \sum_{j=1}^n f(x_j) \right)^2 = \arg\max_{\|f\|_H \leq 1} \frac{1}{n} \sum_{i=1}^n (f(x_i) - \hat{f}(i))^2 = \arg\max_{\|f\|_H \leq 1} \text{Var}(f)
\]

We can write $f = \sum_{i=1}^n \left( \phi(x_i) - \frac{1}{n} \sum_{j=1}^n \phi(x_j) \right) = \sum_{i=1}^n \alpha_i \tilde{\phi}(x_i)$ since any orthogonal component contributes.

Define infinite dimensional analog of covariance: $C = \frac{1}{n} \sum_{i=1}^n \tilde{\phi}(x_i) \otimes \tilde{\phi}(x_i)$, where $\tilde{\phi}(x_i)$ is demeaned features, $\otimes$ is outer product.

with trick $(A \otimes b) c := \langle b, c \rangle_H a$:

\[
C f^T = \frac{1}{n} \sum_{i,j} \tilde{\phi}(x_i) \left( \langle x_j \in E(x_i, x_j) \rangle \right), \quad \tilde{K} = H K H \quad \text{since } \tilde{K} = H^T X^T X H = H^T E H = HEH
\]

project $f_i \alpha = C f_i$ and $\tilde{\phi}(x_i) \Rightarrow n\lambda_i \tilde{\alpha}_i = \tilde{K} \tilde{\alpha}_i \Rightarrow n\lambda_i \alpha = K \alpha$

$\|f\|_H^2 = \alpha^T \tilde{K} \alpha = n \lambda \|\alpha\|^2$, so to make it unit norm, we scale $\alpha \in \alpha / \sqrt{n \lambda}$.

\section{Ridge Regression}

\[
X = [x_1, \ldots, x_n]^T \in \mathbb{R}^{p \times n}, \quad y := [y_1, \ldots, y_n]^T
\]

\[
\alpha^* = \arg\min_{\alpha} \left( \sum_{i=1}^n (y_i - \alpha^T x_i)^2 + \lambda \|\alpha\|^2 \right) = \arg\min_{\alpha \in \mathbb{R}^p} \left( \|y - X^T \alpha\|^2 + \lambda \|\alpha\|^2 \right)
\]

\[
\alpha^* = (XX^T + \lambda I)^{-1} Xy = X(X^T X + \lambda I)^{-1} y
\]

So solution is in form of $\alpha^* = \sum_{i=1}^n \alpha_i^* x_i$.

Perform SVD on $X$: $X = USV^T$, with $U = [u_1, \ldots, u_c]$, $S = [\sigma_1, \sigma_2, \ldots]^T$, $V = [v_1, \ldots, v_n]^T$, $U^T U = UU^T = I_c$, $V^T V = VV^T = I_n$.

\section{Kernel Ridge Regression}

\[
\alpha^* = \arg\min_{\alpha} \left( \sum_{i=1}^n (y_i - \langle \alpha, \phi(x_i) \rangle_H)^2 + \lambda \|\alpha\|_H^2 \right)
\]

\[
\alpha^* = X(K + \lambda I_n)^{-1} y - \sum_{i=1}^n \alpha_i^* \phi(x_i), \quad \alpha = (K + \lambda I_n)^{-1} y
\]

\newpage

\section{Gaussian Kernel}

\[
k(x, x') = \exp\left(-\frac{\|x - x'\|^2}{\sigma^2}\right)
\]

$\sigma$, $k$ less centralized, more smoother. $\|f\|_{H_k}^2 = \sum_j \frac{|\hat{f}_j|^2}{\lambda_j}$, so larger $k$ decays faster, $f$ has fat tail become less penalty $k$!

\section{Representer Theorem}

Let $H$ be RKHS on a set $X$ with kernel $k$. $(x_1, \ldots, x_n) \subset X$ be training points. Let $L : \mathbb{R}^n \to \mathbb{R} \cup \{\infty\}$ be any loss that depends only on the vector of point-wise values:

$L \mapsto (f(x_1), \ldots, f(x_n))$

$\Omega : H_k \to [0, \infty)$ be a non-decreasing function of RKHS norm, e.g., $\Omega(f) = \lambda \|f\|_{H_k}^2$.

Consider the regularized empirical-risk problem:

\[
\min_{f \in H_k} F(f) := L(f(x_1), \ldots, f(x_n)) + \Omega(f)
\]

Any minimizer of $F$ lies in the finite-dimensional span of the kernel sections at the data points:

\[
f(x) = \sum_{i=1}^n \alpha_i k(x, x_i) \quad \text{for some coefficients } (\alpha_1, \ldots, \alpha_n) \in \mathbb{R}^n \text{. If } \Omega \text{ is strictly increasing, } f \text{ is unique}
\]

\section{Kernel Mean Embedding}

We have RKHS $H$ built from kernel $k : X \times X \to \mathbb{R}$ and a Borel probability measure $P$ on $X$.

\textbf{Kernel mean embedding of $P$:}

\[
\mu_P := \mathbb{E}_P[k(\cdot, X)] \in H, \quad \text{such that}
\]

\[
\text{for } \forall f \in H, \quad \mathbb{E}_P[f(X)] = \langle f, \mu_P \rangle_H
\]

Define linear functional $T_P : H \to \mathbb{R}$, $T_P f = \mathbb{E}_P[f(x)]$. With $f(x) = \langle f, \phi(x) \rangle_H$, we have:

\[
|T_P f| = |\mathbb{E}_P[f(X)]| \leq \mathbb{E}_P[|f(X)|] = \mathbb{E}_P[|\langle f, k(\cdot, X) \rangle_H|] \leq \mathbb{E}_P[\|f\|_H \|k(\cdot, X)\|_H]
\]

If $\mathbb{E}_P[\|k(\cdot, X)\|_H] < \infty$, $T_P$ is bounded.

\[
\text{So } \exists! \mu_P \in H \text{ such that } T_P f = \langle f, \mu_P \rangle_H \text{ by Riesz}
\]

Let $f = \phi(x) = k(\cdot, x)$, we have:

\[
\mu_P(x) = \langle \mu_P, \phi(x) \rangle_H = T_P \phi(x) = \mathbb{E}_P[k(x, X)]
\]

$\mu_P$ is just the ordinary mean of those feature vectors. $\mu_P = \mathbb{E}_P[\phi(\theta)]$.

\[
\langle \mu_P, \mu_Q \rangle_H = \mathbb{E}_P \mathbb{E}_Q [\langle k(\cdot, X), k(\cdot, Y) \rangle_H] = \mathbb{E}_P \mathbb{E}_Q [k(X, Y)]
\]

distance between feature means.

\section{Empirical Mean Embedding}

\[
\hat{\mu}_P = \frac{1}{m} \sum_{i=1}^m \phi(X_i), \quad X_i \sim P
\]

\section{Maximum Mean Discrepancy}

\[
\text{MMD}^2(P, Q) = \|\mu_P - \mu_Q\|_H^2 = \mathbb{E}_P[\langle(x, X)\rangle] + \mathbb{E}_Q[\langle(y, y)\rangle]
\]

\[
\text{distance between feature means}
\]

\textbf{empirical MMD:} 

\[
\widehat{\text{MMD}}^2 = \frac{1}{n(n-1)} \sum_{i \neq j} k(x_i, x_j) + \frac{1}{n(n-1)} \sum_{j \neq j'} k(y_j, y_{j'}) - \frac{2}{n^2} \sum_i k(x_i, y_j)
\]

\newpage

\section{MMD as an Integral Probability Metric}

Find a ``well learned function'' $f(x)$ to maximize $\mathbb{E}_P[\mathbb{E}[f(X)] - \mathbb{E}_Q[f(D)]]$.

\[
\text{MMD}(P, Q; F) = \|\mu_P - \mu_Q\|_H^2 = \|T_P - T_Q\|_{F'} = \sup_{\|f\|_H \leq 1} |\langle T_P - T_Q f \rangle| = \sup_{\|f\|_H \leq 1} |\mathbb{E}_P[f(X)] - \mathbb{E}_Q[f(X)]|
\]

$F$ is unit ball in RKHS $F$.

$\text{MMD}(P, Q; F) = 0 \Leftrightarrow P = Q$.

\section{Witness Function}

\textbf{Witness function:} $f^* \sim \mu_P - \mu_Q$.

\[
f^*(v) = \langle f^*, \phi(v) \rangle_H \propto \langle \mu_P - \mu_Q, \phi(v) \rangle_H = \frac{1}{n} \sum_{i=1}^n k(x_i, v) - \frac{1}{n} \sum_{j=1}^n k(y_j, v)
\]

\section{Hypothesis Testing on MMD}

Null hypothesis $H_0$: when $P = Q$ should see $\widehat{\text{MMD}}^2$ close to 0. When $P \neq Q$, statistic is asymptotically normal:

\[
\frac{\widehat{\text{MMD}}^2 - \text{MMD}^2(P, Q)}{\sqrt{V_n(p, Q)}} \xrightarrow{D} \mathcal{N}(0, 1)
\]

when $P = Q$, statistic has asymptotic distribution $n \widehat{\text{MMD}}^2 \sim \sum_{\ell=1}^{\infty} \lambda_\ell [\mathbb{Z}_\ell^2 - 2]$

$\mathbb{Z}_\ell \sim \mathcal{N}(0, 2)$ i.i.d.

We can permute samples from $P$ and $Q$ to simulate $P = Q$ and get $\alpha$ (false positive rate bound).

\section{Maximum Mean Embedding via Fourier Series}

\[
\mu_P(x) = \langle \mu_P, k(\cdot, x) \rangle_H = \mathbb{E}_P[\mathbb{E}[k(t - x)]] = \int_X k(t - x) d P(t)
\]

\[
\text{MMD}(P, Q; H(\mathcal{K})) = \|\mu_P - \mu_Q\|_H^2 = \left\| \sum_{\ell=0}^{\infty} [\phi_{\rho,\ell} - \phi_{Q,\ell}] \mathcal{K}_\ell^2 \exp(i \omega_\ell x) \right\|_H,
\]

\textbf{MMD}$^{\text{c}}(P, Q; H^f) = \sum_{\ell=0}^{\infty} |\varphi_{\rho,\ell} - \varphi_{Q,\ell}|^2 \mathcal{K}_\ell^2 = \sum_{\ell=0}^{\infty} |\varphi_{\rho,\ell} - \varphi_{Q,\ell}|^2 \mathcal{K}_\ell$

$\mathcal{K}(x, y) = k(x - y) = \sum_{\ell=0}^{\infty} \mathcal{K}_\ell e^{i \omega_\ell (x - y)}$, $\mathcal{K}_\ell > 0$

$\psi_{\rho,\ell} := \int e^{i \ell t} d P(t)$ (characteristic function).

$\mu_P(x) = \int_X k(t - x) d P(x) = \int_X \sum_{\ell} \mathcal{K}_\ell e^{i \omega_\ell(t-x)} d P(t) = \sum_\ell \mathcal{K}_\ell e^{i \omega_\ell x} \int_X e^{-i \omega_\ell t} d P(t) = \sum_\ell \mathcal{K}_\ell e^{i \omega_\ell x} \psi_{\rho,\ell}$

(we can use $e^{i \omega x}$ since $\mathcal{K}_\ell = \mathcal{K}_{-\ell}$)

$|\psi_{\rho,\ell} - \hat{\varphi}_\ell|^2 = |\varphi_{\rho,\ell} - \varphi_{\ell+m}|^2$ since $[\hat{\mathcal{K}}_\ell^2 = 2 \tilde{\mathcal{K}}$.

\section{Kernel is Characteristic}

A kernel is characteristic if the mean-embedding map $P \mapsto \mu_P := \int k(\cdot, x) d P(x) \in H_k$ is injective, i.e., $\mu_P = \mu_Q \Rightarrow P = Q$.

For translation-invariant kernel $\mathcal{K}(x, y) = k(x - y)$, $\mu_P = \int k(\cdot - x) d P(x) \in H_k$.

$\mu_P$ is convolution of $\mathcal{K} * P$, by convolution theorem, $\widehat{\mathcal{K} * P}(\omega) = \hat{\mathcal{K}}(\omega) \hat{P}(\omega)$, where $\mathcal{K}(t) = (\int e^{-i \omega x t} d \Lambda(\omega)$, $\hat{P}(\omega) = (\int e^{-i \omega x} d P(x)$.

so $\mu_P = \mu_Q \Leftrightarrow \hat{\mathcal{K}}(\omega) [\hat{P}(\omega) - \hat{Q}(\omega)] = 0$ $\forall \omega$.

So full-support spectrum $\hat{\mathcal{K}}(\omega) \neq 0$ $\forall \omega$ forces $\hat{P} = \hat{Q}$ everywhere, so $P = Q$.

In continuous case, $\mathcal{K}(t) = \int e^{i \omega x t} d \Lambda(\omega)$, $\hat{P}(\omega) = \int e^{-i \omega x} d \Lambda(\omega)$, $\hat{P}(\omega) = \frac{d \Lambda(\omega)}{d \omega}$.

In periodic domain case, spectrum is a sum $\{\lambda_\ell\}$ (continuous)

so kernel is characteristic $\Leftrightarrow$ \{\}

\newpage

\section{Fourier Representation of MMD on $\mathbb{R}^d$}

\[
\text{MMD}^2(P, Q; F) = \int |\hat{P}(\omega) - \hat{Q}(\omega)|^2 d \Lambda(\omega)
\]

\section{Properties of MMD and Kernels}

\begin{enumerate}
    \item \textbf{MMD decay:} with increasing perturbation frequency as $\mathcal{K}$ decays as $\ell$ grows.
    
    \item \textbf{Translation invariant kernel:} $k$ is characteristic for probability measures on $\mathbb{R}^d \Leftrightarrow \text{supp}(\Lambda) = \mathbb{R}^d$ (support zero on at most a countable set).
    
    \item \textbf{Any continuous compactly supported kernel:} is characteristic (since Fourier spectrum $\Lambda(\omega)$ cannot be 0 on an interval).
    
    \item \textbf{Universal RKHS:} $k(x, x')$ continuous, $X$ compact, and $F$ dense in $C(X)$ with respect to $\|\cdot\|_\infty$.
    
    If $X$ is compact, continuous $k(\cdot, x)$ creates compact-image.
    
    If $X$ is not compact, $\|k(\cdot, x)\|_F < \infty$.
    
    $F$ is dense in $C(X) \Rightarrow$ under $\|\cdot\|_\infty \Leftrightarrow \forall g \in C(X)$, $\forall \varepsilon > 0$, $\exists f \in F$ such that $\|f - g\|_\infty < \varepsilon$.
    
    \item \textbf{Universal $F \Rightarrow$ kernel is characteristic}
\end{enumerate}

\section{Test Power and Kernel Selection}

Test power depends on $k(x, x')$, $P$ and $Q$ (and $n$), with characteristic kernel, MMD has power $\to 1$ as $n \to \infty$ for any fixed problem. But could be terrible for reasonable $n$.

\[
\mathbb{P}[n \widehat{\text{MMD}}^2 \geq \xi] \Rightarrow \mathbb{E}\left( \frac{\widehat{\text{MMD}}(P, Q) - c_n}{\sqrt{V_n(p, Q)}} \right), \quad \text{second term is } O(\hat{n})
\]

to maximize test power, maximize $\frac{\widehat{\text{MMD}}(P, Q)}{\sqrt{V_n(p, Q)}}$ for fix $n$.

\textbf{Split data first} to choose a kernel $k$ maximizing the term and use the $k$ to perform MMD test on rest of data.

\section{Type I and Type II Errors}

\begin{enumerate}
    \item \textbf{Type I error:} false positive $\to$ reject a true $H_0$
    \item \textbf{Type II error:} false negative $\to$ accept a false $H_0$
\end{enumerate}

\section{MMD as a Dependence Measure}

\[
\text{MMD}[\rho_{XY}, P_X P_Y, H_k] = \text{Simulate } Q \text{ with pairs } (x_i, y_j)
\]

\section{Constrained Covariance}

\[
\text{COCO}(\rho_{XY}) = \sup_{\|f\|_F=1, \|g\|_G=1} \text{cov}[f(x)g(y)], \quad f \in F, g \in G, \text{two spaces}
\]

$\text{COCO} = $ max singular value of feature covariance $C_(\rho(f,g))$.

\section{Hilbert-Schmidt Operators}

Let $F$ and $G$ be separable Hilbert spaces. $\{e_i\}_{i \in J}$ are orthonormal basis of $F$, $\{f_j\}_{j \in I}$ for $G$.

Define two compact linear operators $L : G \to F$ and $M: G \to F$:

\[
\|L\|_{HS}^2 = \sum_{j \in J} \|L f_j\|_F^2 = \sum_{i \in J} \sum_{j \in I} |\langle L f_j, e_i \rangle_F |^2
\]

Hilbert-Schmidt operators form a Hilbert space $HS(G, F)$ with inner-product:

\[
\langle L, M \rangle_{HS} = \sum_{j \in J} \langle L f_j, M f_j \rangle_F = \sum_{i \in I} \sum_{j \in J} \langle L f_j, e_i \rangle_F \langle M f_j, e_i \rangle_F
\]

\section{Finite Dimensions and Frobenius Norm}

In finite dimensions, if $L$ and $M$ are matrices, the HS inner product reduces to the Frobenius inner product:
\[
\langle L, M \rangle_{HS} = \text{tr}(L^T M), \quad \|L\|_{HS}^2 = \sqrt{\text{tr}(L^T L)} = \sum_{i,j} |L_{ij}|^2, \quad \hat{\mu}_i := \frac{1}{n} \sum_{i=1}^{n} \phi(x_i)
\]

Hilbert-Schmidt inner product measures similarity of $L$ and $M$ on basis functions.

If $L$ and $M$ transform OMS in a similar way, their HS inner product is large.

\section{Rank-One Operator}

A rank-one operator $T: G \to\mathbb{R}F$ can be written as:
\[
T = a \otimes b, \quad a \in F, b \in G
\]

For $\psi \in G$, $(a \otimes b)\psi = \langle b, \psi \rangle_G a$.

The image of $T$ is the one-dimensional subspace $\text{span}\{a\}$, hence rank-one:

\[
\|a \otimes b\|_{HS}^2 = \sum_{j \in J} \|(a \otimes b)f_j\|_F^2 = \sum_{j \in J} |\langle b, f_j \rangle_G|^2 a\|_F = \|a\|_F^2 \sum_{j \in J} \|b, f_j \rangle_G\|_F^2 = \|a\|_F^2 \|b\|_G^2 < \infty
\]

$a \otimes b$ is a Hilbert-Schmidt operator.

\section{Expected Hilbert-Schmidt Inner Product Functional}

Define linear functionals operator by expected Hilbert-Schmidt inner product functional:

\[
T_{xy} : HS(G, F) \to \mathbb{R}
\]

\[
T_{xy}(A) = \mathbb{E}_{xy} \left[ \left\langle \phi(x) \otimes \psi(y), A \right\rangle_{HS} \right]
\]

It is a bounded operator if:
\[
\mathbb{E}_{xy} \left\| [\phi(x) \otimes \psi(y)] \|_{HS} \right] = \mathbb{E}_{xy} \left[ [\phi(x) \otimes \psi(y)]^* \right] < \infty
\]

then $T_{xy} \in HS^* $ (a continuous dual of HS)

and by Riesz representation theorem: $\exists ! \hat{C}_{xy} \in HS(G,F)$ such that:

\[
T_{xy}(A) = \langle \hat{C}_{xy}, A \rangle_{HS}
\]

\section{Rank-One Covariance Operator}

We choose $A = f \otimes g$. Then:
\[
\langle \hat{C}_{xy}, f \otimes g \rangle_{HS} = \langle f, \hat{C}_{xy} g \rangle_F
= \mathbb{E}_{xy} [\langle \phi(x), f \rangle_F \langle \psi(y), g \rangle_F]
\]

Define by expected operator:
\[
T_{xy}: HS(G, F) \to \mathbb{R}, \quad T_{xy}(A) = \mathbb{E}_{xy} [\langle \phi(x) \otimes \psi(y), A \rangle_{HS}]
\]

\section{Empirical Covariance Operator}

\[
\hat{C}_{xy} = \frac{1}{n} \sum_{i=1}^{n} \phi(x_i) \otimes \psi(y_i) - \mu_x \otimes \mu_y, \quad \text{where} \quad \mu_x := \frac{1}{n} \sum_{i=1}^{n} \phi(x_i)
\]

\[
\hat{C}_{xy} = \frac{1}{n} X^T Y, \quad X = [\phi(x_1), \ldots, \phi(x_n)]^T, \quad Y = [\psi(y_1), \ldots, \psi(y_n)]^T
\]

\[
\tilde{C}_{xy} = HR H, \quad \tilde{C}_{xy} = HLH
\]

\[
\text{COCO}(P_{xy}) = \sup_{\|f\|=1, \|g\|=1} \text{cov}[f(x)g(y)]
\]

We can do maximize $\langle f, \hat{C}_{xy} g \rangle_F$ subject to $\|f\|_F = 1, \|g\|_G = 1$. Witness function should be $f = \frac{1}{\sqrt{n}} \sum_i a_i [\phi(x_i) - \mu_x]$ and $g = Y H P$.

$f \hat{C}_{xy} g = \frac{1}{n} \alpha^T K L \beta$, and $\|f\|_F^2 = \alpha^T K \alpha$, $\|g\|_G^2 = \beta^T L \beta$.

\section{Associated Lagrangian}

\[
\text{COCO}(P_{xy}) = \sup_{\|f\|=1, \|g\|_G=1} \text{cov}[f(x)g(y)]
\]

With Lagrangian formula with multipliers $\gamma \in \mathbb{R}^m$:

\[
\mathcal{L}(P, Q; F, g) = - f^T \hat{C}_{xy} g + \frac{1}{2}[\|f\|_F^2 - 1) + \frac{1}{2}c\|g\|_G^2 - 1]
\]

where $\alpha_0 = \sup (\mathbb{E}_{xy} [\phi(x) \otimes \psi(y)] \|_{HS})$, $\alpha = r$.

\section{Hilbert-Schmidt Independence Criterion (HSIC)}

\[
\text{HSIC}^2(\mathcal{P}_X, \mathcal{P}_Y; F, G) = \sum_{i=1}^{\infty} \delta_i^2, \quad \delta_i := \text{CoCD}(\mathcal{P}_{XY}; F, G)
\]

HSIC is MMD with product kernel

\[
\text{HSIC}(P_{XY}; F, G) = \text{MMD}(\mathcal{P}_X, P_X P_Y; H_k) \quad \text{where } k((x,y), (x',y')) = k(x, x') l(y, y')
\]

\[
\text{HSC}^2(P_{XY}; F, G) := \| \hat{C}_{XY} - \mu_X \otimes \mu_Y \|_{HS}
\]

\[
= \langle \hat{C}_{XY}, \hat{C}_{XY} \rangle_{HS} + \langle \mu_X \otimes \mu_Y, \mu_X \otimes \mu_Y \rangle_{HS} - 2 \langle \hat{C}_{XY}, \mu_X \otimes \mu_Y \rangle_{HS}
\]

\[
\| \hat{C}_{XY} \|_{HS}^2 = \mathbb{E}_X [k(x, x')] \mathbb{E}_Y [\ell(y, y')] = D \cdot B
\]

\[
\| \mu_X \otimes \mu_Y \|_{HS}^2 = \| \mu_X \|_F^2 \| \mu_Y \|_G^2 = \mathbb{E}_{XY} [k(x,x')] \mathbb{E}_Y[[\ell(y, y')]] =: B
\]

\[
\langle \hat{C}_{XY}, \mu_X \otimes \mu_Y \rangle_{HS} = \mathbb{E}_{XY} [[\phi(x) \otimes \psi(y), \mu_X \otimes \mu_Y \rangle_{HS}] = B
\]

So $\text{HSIC}^2 = A - 2B + D$

\section{Unbiased and Biased HSIC Estimators}

\textbf{Unbiased } $\hat{A} := \frac{1}{n(n-3)} \sum_{i=1}^{n} \sum_{j \neq i, j'} k(x_i, x_j) \ell(y_i, y_j)$.

\textbf{Biased } $\hat{A}_b := \langle \frac{1}{n} \sum_{i \phi(x_i) \otimes \psi(y_i), \frac{1}{n}\sum_{i \phi(x_i) \otimes \psi(y_i)} \rangle_{HS}$

$= \frac{1}{n^2} \sum_{i,j} k(x_i, x_j) \ell(y_i, y_j) = \frac{1}{n^2} \text{tr}(KL)$.

Expectation of bias is of $O(\frac{1}{n})$.

\section{Empirical HSIC}

\[
\text{HSIC}^2 = \frac{1}{n^2} \text{tr}(KLH), \quad \text{HSIC}^2(Q) = \sum_{i=1}^{\infty} \lambda_l z_\ell^2, \quad z_\ell \sim \mathcal{N}(0,1) \text{ i.i.d}.
\]

\section{Asymptotics of HSIC}

when $P_{XY} = P_X P_Y P^*$: $n \text{HSIC}^2 \Rightarrow \sum_{\ell=1}^{\infty} \lambda_\ell z_\ell^2$, $z_\ell \sim \mathcal{N}(0,1)$ i.i.d.

\section{Null Distribution}

Given $P_{XY} = P_X P_Y$, what is threshold $C_\alpha$ such that $\mathbb{P}(\text{HSIC}^2 > C_\alpha) < \alpha$?

\textbf{Under null distribution via permutation:} permute $y$ and calculate HSIC many times to get empirical CDF.

\textbf{Under null distribution via moment matching:}

\[
n \text{HSIC}^2(\hat{Z}) \sim \frac{\alpha^2}{n^2} T(\alpha), \quad \alpha = \frac{\mathbb{E}(\text{HSIC}^2)^2}{\text{Var}(\text{HSIC}^2)}, \quad \beta = \frac{n\mathbb{E}(\text{HSIC}^2)}{C_\alpha}
\]

\section{V-structure Discrepancy}

\textbf{CONSISTS CI test:} $H_0 : X \perp Y | Z$ \quad \textbf{Factorization test:} $H_0: (X,Y) \perp Z | T V(Y,Z) \perp X$.

\section{Lancaster Interaction Measure}

Lancaster interaction measure of $(X_1, \ldots, X_D) \wedge P$ is a signed measure $\Delta P$ that vanishes whenever $P$ can be factorised non-trivially.

\[
D=2: \Delta P = P_{XY} - P_X P_Y, \quad D=3: \Delta P = P_{X_1 X_2 X_3} - P_{X_1} P_{X_2 X_3} - P_{X_2} P_{X_1 X_3} - P_{X_3} P_{X_1 X_2} + 2 P_{X_1} P_{X_2} P_{X_3}
\]

\[
(X, Y) \perp Z \Rightarrow \Delta P = 0
\]

but $\Delta_1 P = 0 \Rightarrow (X, Y) \perp Z \vee (X, Z) \perp Y \vee (Y, Z) \perp X$. It is only useful when $\Delta P \neq D$.

\section{HSIC on Feature Selection}

Backwards elimination of irrelevance features to maximize HSIC.

\section{Kernel Stein Discrepancy (KSD)}

\[
\text{KSD}(\pi) := \sup_{\|g\|_{\ell_\infty} \leq 1} \mathbb{E}_\pi[[\nabla_x \log \pi(x^*)]
\]

\[
= \sup_{\|g\|_{\ell_\infty} \leq 1} \mathbb{E}_\pi[\nabla_x \log p(x)]
\]

\section{Langevin Stein Operator}

\[
\mathcal{A}_p^f(x) = \frac{1}{p(x)} \frac{d}{dx} [f(x) p(x)] = \nabla_x \log p(x) \cdot f(x) + \nabla_x f(x)
\]

\[
\mathbb{E}_p [\mathcal{A}_p^f] = 0 \quad \text{subject to} \quad [f(x) p(x)]_\infty = 0
\]

Here we do not need to normalize $p(x)$.

\section{Kernel Stein Discrepancy (KSD)}

\[
\text{KSD}_p(\mathcal{Q}) = \sup_{\|g\|_F \leq 1} \mathbb{E}_Q[\mathcal{A}_p g]
\]

\[
= \sup_{\|g\|_F \leq 1} \mathbb{E}_Q[g].
\]

\section{Reproducing Property for the Derivative}

\[
\frac{d}{dx} f(x) = \left\langle f, \frac{d}{dx} k(x, x') \right\rangle_F, \quad \left\langle \frac{d}{dx} \phi(x), \psi(x) \right\rangle_F = \frac{d}{dx} k(x, x)
\]

\section{Derivatives of Functions}

\[
[\mathcal{A}_p f](x) = \left( \frac{d}{dx} \log p(x) \right) f(x) + \frac{d}{dx} f(x) = \left\langle f, \frac{d}{dx} k(x, x') \right\rangle_F
\]

\[
z_\pi(x) = \left( \frac{d}{dx} \log p(x) \right) (\psi(x)) + \frac{d}{dx} \psi(x)
\]

\section{Expected Stein Operator}

\[
T_{\rho, p^*} = \mathbb{E}_Q [\mathcal{A}_p f(x)] = \left\langle f, \mathbb{E}_Q \left[ \frac{d}{dx} \log p(x) \right] \psi(x) + \frac{d}{dx} \psi(x) \right\rangle_F
\]

\[
( T_{\rho, p} : H \to \mathbb{R} )
\]

\[
\left| \mathbb{E}_Q [\langle f, z_x \rangle_F] \right| \leq \mathbb{E}_Q [|\langle f, z_x \rangle_F|] \leq \| f \|_F \mathbb{E}_Q [|| z_x ||_F]
\]

\[
|| z_x ||_F^2 = C + \left( \frac{d}{dx} \log p(x) \right)^2 c, \quad \text{for boundedness we need to have } \mathbb{E}_Q \left[ \left( \frac{d}{dx} \log p(x) \right)^2 \right] < \infty
\]

So that $\mathbb{E}_Q [|| z_x ||_F] = \mathbb{E}_Q \left[ C + \left( \frac{d}{dx} \log p(x) \right)^2 c \right] \leq \sqrt{\mathbb{E}_Q \left[ c + \frac{d}{dx} \log p(x) \right]^2 c} < \infty$.

If data is heavy-tailed and model has right tail, KSD may not be defined.

So by Riesz, $T_\rho, p^* = \mathbb{E}_Q [\mathcal{A}_p f(x)] = \langle f, z_x \rangle_F$.

\section{KSD Properties}

\[
\text{KSD} = \sup_{\|g\|_F \leq 1} \mathbb{E}_Q[\mathcal{A}_p g] = \sup_{\|g\|_F \leq 1} \mathbb{E}_Q [g], \quad \mathbb{E}_Q[z_x] \sim \mathcal{F}
\]

witness function should be $g \propto \mathbb{E}_Q[z_x]$.

\section{Squared KSD}

KSD can be written as $\text{KSD}^2(\mathcal{Q}, P) = \mathbb{E}_{x, x' \sim Q}[k_p(x, x')]$, where $k_p(x, \ell) = \mathcal{A}_p^x \mathcal{A}_p^y k(x, x')$, and $\mathcal{A}_p^{(1)}, \mathcal{A}_p^{(2)}$ act on the first and second arguments of $k$.

\[
\mathcal{K}_p(x, y) = S_p(x)^T S_p(x) k(x, y) + S_p(x)^T k_2(x, y) + k_1(x, y) k(x, y) + \sigma[k_{\alpha_1}(x, y)]
\]

\[
S_p(x) = \frac{\nabla p(x)}{p(x)}, \quad k_1(x, x') = \nabla_x k(x, x'), \quad k_2(x, x') = \nabla_y k(x, x'), \quad h_p(x, x') = \nabla_x \nabla_y k(x, x')
\]

\section{Estimating KSD}

\[
\text{KSD}(Q, P; F) = \frac{1}{n(n-1)} \sum_{i=1}^{n} \sum_{j \neq i} \delta_i \delta_j h_p(z_i, z_j), \quad \text{where } \{k_i\}_{i=1}^n \text{ i.i.d., } \mathbb{E}(\delta_i) = 0 \text{ and } \mathbb{E}(\delta_i^2) = 1
\]

\[
\sqrt{n} \left( \text{KSD}_p(\hat{Q}) - \text{KSD}_p(Q) \right) \xrightarrow{d} \mathcal{N}(0, \delta_p^2) \quad \text{where } \delta_p^2 = 4\text{Var}(\mathbb{E}_x[\mathcal{A}_p \ell(y, x')])
\]

For latent variable model testing $\hat{p}(x) = \int p(x|z) p(z) dz$, $p(x)$ is usually not tractable.

we use $S_p(x) = \frac{\nabla p(x)}{p(x)} = \frac{\nabla p(x|z)}{p(x)} = \int \frac{\nabla p(x|z)}{p(x|z)} p(z|x) \mathrm{d}z = \mathbb{E}_{z|x} [S_p(x|z)]$.

\section{Comparing Models via Samples}

\textbf{Compare sample from Q and a model P:} usually can't compute $\mathbb{E}_P$ in closed form.

\textbf{Compare a sample and a model:} KSD.

\section{Hyperplane in $\mathbb{R}^D$}

For a hyperplane $\{x \in \mathbb{R}^D : w^T x + b = 0\}$, $w$ is perpendicular to the hyperplane.

\section{Distance of Point to Hyperplane}

Distance of a point $x$ to the hyperplane is $\frac{|w^T x + b|}{\|w\|}$.

\section{Generic Optimization Problem}

Generic optimization problem on $X \in \mathbb{R}^D$:

\[
\text{minimize} \quad f_0(x)
\]

\[
\text{subject to} \quad f_i(x) \leq 0, \quad i = 1, \ldots, m
\]

\[
h_i(x) = 0, \quad i = 1, \ldots, p
\]

\textbf{Lagrangian} $L$ is a lower bound on the original problem. 

\[
L(x, \lambda, \nu) := f_0(x) + \sum_{i=1}^{m} \lambda_i f_i(x) + \sum_{i=1}^{p} \nu_i h_i(x)
\]

\[
I_-(w) = \begin{cases} 0 & w \leq 0 \\ \infty & w > 0 \end{cases}, \quad I_0(w) = \begin{cases} 0 & w = 0 \\ \infty & w \neq 0 \end{cases}
\]

To ensure a lower bound, $\lambda \geq 0$.

\[
\inf_x g(x, \lambda, \nu) = \inf_x L(x, \lambda, \nu) \leq p^*
\]

when $\lambda \geq 0$ and $f_i(x) \leq 0$, Lagrange penalizes out of bounds of $f_i$ and reward is the bound as $\lambda_i f(x) > 0$ when out of bound while $\lambda_i f(x) \leq 0$ when in the bound.

\section{Lagrange Dual Problem}

\textbf{Lagrange dual problem:} maximize $g(\lambda, \nu)$ subject to $\lambda \geq 0$.

\textbf{Dual function is a pointwise infimum} of affine functions of $(\lambda, \nu)$ with different $x$, hence is concave in $(\lambda, \nu)$ with convex constraint set $\lambda \geq 0$.

\[
d^* = \max g(\lambda, \nu) \leq p^*
\]

Strong duality means $d^* = p^*$.

\section{Strong Duality Condition}

\textbf{Strong duality} holds if:
\begin{enumerate}
    \item the primal problem is convex, i.e., $f_0, \ldots, f_m$ are convex minimize $f_o(x)$ subject to $f_i(x) \leq 0$ $i = 1, \ldots, n$, $Ax = b$.
    
    \item \textbf{Slater's condition holds:} there exists some strictly feasible point $x^* \in \text{relint}(\mathcal{D})$ such that $f_i(x^*) < 0$, $i = 1, \ldots, m$, $Ax^* = b$.
    
    \item In affine $f_i$ case, we allow $f_i(x^*) \leq 0$. (Only need strict inequality over non-affine constraints).
\end{enumerate}

\section{KKT Conditions}

Under strong duality, $x^* \in \arg\max_x \min_y L(x, y)$ and $y^*$ are primal and dual solutions if:

\begin{enumerate}
    \item \textbf{Stationarity:} $0 \in \partial f_0(x^*) + \sum_i \lambda_i^* \partial f_i(x^*) + \sum_j \nu_j^* h_j(x^*)$
    
    \item \textbf{Complementary slackness:} $\lambda_i^* f_i(x^*) = 0$, $i = 1, \ldots, m$.
    
    \item \textbf{Primal feasibility:} $f_i(x^*) \leq 0$, $i = 1, \ldots, m$, $h_i(x^*) = 0$, $i = 1, \ldots, p$.
    
    \item \textbf{Dual feasibility:} $\lambda_i^* \geq 0$, $i = 1, \ldots, m$.
\end{enumerate}

\section{Complementary Slackness}

This is a necessary condition for optimality:

\[
\lambda_i^* f_i(x^*) = 0, \quad \text{which means} \quad \lambda_i^* > 0 \Rightarrow f_i(x^*) = 0
\]

\[
f_i(x^*) < 0 \Rightarrow \lambda_i^* = 0
\]

In other $f_i$ case, we allow $f_i(x^*) \leq 0$. (Only need strict inequality over non-affine constraints).

\end{document}
